\par
\SourceLeaf{490}%
Ich gebe Ew. — — unterthänigst anheim, diese Gründe zu prüfen und darüber zu entscheiden, und füge nur noch hinzu, dass freilich diese Abwesenheit mich in Beziehung auf die zwei Collegia, welche ich in diesem Sommer zu halten habe, mich etwas beengen würde, dass ich jedoch bei der kleinen Anzahl der Zuhörer mit diesen ein Arrangement treffen könnte, die ausfallenden Stunden theils vorher theils nachher einzubringen. Und um die Zeit der Abwesenheit nach Möglichkeit zu beschränken, würde ich mit dem Prof. \textsc{Schumacher} (welchen ich jeden Tag hier zurückerwarte) vorläufig mündliche Abrede nehmen, damit er in Lauenburg vorher alle Vorbereitungen treffen und die Beobachtungen unmittelbar nach meiner Ankunft ihren Anfang nehmen könnten. Da in dieser Jahreszeit das Wetter den Beobachtungen nicht ungünstig zu sein pflegt, so wäre zu erwarten, dass die Zeit meiner Abwesenheit nicht länger zu sein brauchte, als etwa die Badereisen, welche manche meiner Collegen um dieselbe Zeit wohl machen.

Ich benutze diese Gelegenheit, um Ew. — — noch anzuzeigen, dass ich von dem Mechanikus \textsc{Rumpf} verfertigte Maschine zur Umlegung des \textsc{Reichenbach}schen Mittagsfernrohrs nunmehr in diesen Tagen angeschlagen wird, und dass dann sofort auch an diesem Instrumente die regelmässigen Beobachtungen ihren Anfang nehmen werden.

— — —

Göttingen den\hspace{3em}1819.
\begin{flushright}C. F. \textsc{Gauss}.\end{flushright}

\par\bigskip
\pdfbookmark[0]{Anzeige betreffend die Fortsetzung der Dänischen Gradmessungen. 1820.}{anzeige-gradmessung-1820}
\begin{center}\large ANZEIGE DES HOFR. GAUSS, BETREFFEND DIE FORTSETZUNG DER DÄNISCHEN GRADMESSUNGEN DURCH DAS KÖNIGREICH HANNOVER.\end{center}
\par
— — —

Zenithsectoren wurden bisher auch für die ersten Sternwarten als ein wesentliches Bedürfniss angesehen. Allein gegenwärtig machen die neuconstruirten Meridiankreise, namentlich die \textsc{Reichenbach}schen, jene Instrumente gewissermassen entbehrlich. Wenigstens leistet, nach meinen bisherigen Erfahrungen der hiesige seit Februar d. J. im Gebrauch befindliche Meridiankreis alles was ein Zenithsector leisten könnte, ebenso vollkommen. Es könnte daher überflüssig scheinen, bei der Gradmessung noch an einen Zenithsector zu denken, wenn nicht folgende zwei wichtige Umstände noch zu berücksichtigen wären.

1) Der Meridiankreis erfordert eine sehr solide Aufstellung, die ihm nur auf einer eigentlichen Sternwarte gegeben werden kann, und ist überhaupt als ein fixes, nicht als ein transportables Instrument zu betrachten. Obgleich nun aber die \textit{wichtigsten} Beobachtungen, wozu ein Zenithsector oder der Meridiankreis im Bezug auf die Gradmessung verwandt werden, die auf der hiesigen Sternwarte anzustellenden sein werden, die in sofern als südlicher Endpunkt betrachtet werden muss, und obgleich diesseitige Beobachtungen der Art am nordlichen Ende deswegen überflüssig sind, oder scheinen können, weil dazu schon die Dänischer Seits in Lauenburg gemachten dienen können, so könnte sich doch beim Verfolg der Arbeit die Nothwendigkeit zeigen, und auf alle Fälle wird es wissenschaftlich interessant sein, dass ähnliche Beobachtungen auf einem Zwischenpunkte des Königreichs Hannover, z. B. in Celle oder Hannover, angestellt würden, was nur mit Hülfe eines transportabeln Zenithsectors geschehen könnte.

2) Da der Gebrauch, der bei einer Gradmessung von den hier in Rede stehenden Beobachtungen gemacht wird, ganz auf der \textit{Vergleichung} der an den verschiedenen Hauptpunkten gemachten Beobachtungen unter \textit{einander} beruht, und also demnächst die Göttinger und vielleicht Celler oder andere diesseitige Beobachtungen mit denen, die Dänischer Seits in Lauenburg, Lysabbel (auf der Insel Alsen) und
